\section*{Chapitre \ref{ch:rs:price-and-volume-features} --- Caractéristiques de prix et de volume}

\begin{solution}{exo:rs:price-and-volume-features:1}
$\ln(51.20/49.80) = 0.0277$~; $(0.0277)^2/(4\ln2) = 0.000277$, soit une volatilité quotidienne de 1.67~\%.
\end{solution}

\begin{solution}{exo:rs:price-and-volume-features:2}
Le momentum 3-1 exclut le dernier mois~: $1.02 \times 0.99 \times 1.04 - 1 = 5.02~\%$. Le rendement sur quatre mois vaut $1.050 \times 1.03 - 1 = 8.17~\%$.
\end{solution}

\begin{solution}{exo:rs:price-and-volume-features:3}
$1.2~\%/40 = 0.03~\%$ par million de dollars~; $2.5~\%/2 = 1.25~\%$ par million de dollars, soit environ 40 fois moins liquide selon cette mesure.
\end{solution}

\begin{solution}{exo:rs:price-and-volume-features:4}
$w_j = (5 - j)/6 - (2 - j)^+/3$~: $w_0 = 1/6$, $w_1 = 1/3$, $w_2 = 1/2$, $w_3 = 1/3$, $w_4 = 1/6$~; la somme vaut $1.5 = (6 - 3)/2$.
\end{solution}

\begin{solution}{exo:rs:price-and-volume-features:5}
La variance d'une moyenne diminue avec le nombre de jours, si bien que $60/5.2 = 11.5$ jours d'estimations de Parkinson égalent soixante
carrés de rendements.
\end{solution}

\begin{solution}{exo:rs:price-and-volume-features:6}
$(\sqrt{11} \times 0.03 + 1 \times (-0.03))/\sqrt{12} = (0.0995 - 0.03)/3.464 = 0.020$, soit un tiers de moins que le 0.03 de la partie 12-1
seule~: avec un retournement mensuel, la \omterm{def:rs:price-and-volume-features:return}{période d'exclusion} vaut un tiers de l'IC.
\end{solution}

\begin{solution}{exo:rs:price-and-volume-features:7}
Avec 30 pas, les \omterm{def:rs:price-and-volume-features:range}{estimateurs de volatilité par l'étendue} affichent des valeurs inférieures de 22--30~\% (Parkinson 0.78, Garman--Klass et Rogers--Satchell 0.70), contre 8--10~\% avec 390~:
les vrais plus haut et plus bas tombent entre les points échantillonnés plus souvent, et d'un écart d'autant plus grand que ces points sont plus espacés. Les barres construites à partir de peu de transactions
(actions peu liquides) ont le même problème.
\end{solution}

\begin{solution}{exo:rs:price-and-volume-features:8}
Une fenêtre centrée sur aujourd'hui inclut le volume des dix jours suivants, qui n'est pas connu à la clôture~: la caractéristique regarde vers l'avant. Le test de fuite
d'information le détecterait. Utiliser les 21 jours précédents, à l'exclusion d'aujourd'hui.
\end{solution}

\begin{solution}{pb:rs:price-and-volume-features:1}
\textbf{1.} Pour un mouvement brownien observé en continu, l'espérance du carré de l'étendue de son logarithme sur une journée vaut $4\ln2\,\sigma^2$.
\textbf{2.} Parkinson. \textbf{3.} Rogers--Satchell~: chacun de ses deux produits associe un plus haut (ou un plus bas) mesuré depuis la clôture et depuis
l'ouverture, si bien qu'une dérive intervient avec des signes opposés et s'annule. \textbf{4.} La variance overnight~; $k = 0.34/(1.34 + 21/19) = 0.139$.
\textbf{5.} 5.2, 7.6 et 5.7. \textbf{6.} 7.5. \textbf{7.} Environ 7.6. \textbf{8.} 0.92, 0.90, 0.90 et 0.91~: l'étendue observée à 390
points est plus étroite que celle observée en continu. \textbf{9.} Le close-to-close passe à 1.25 (il mesure $\sigma^2 + \mu^2$), Parkinson à
1.02, Rogers--Satchell reste à 0.90. \textbf{10.} Les \omterm{def:rs:price-and-volume-features:range}{estimateurs de volatilité par l'étendue} sur un seul jour tombent à 0.63--0.64~; Yang--Zhang reste à 0.93.
\textbf{11.} Le rendement close-to-close couvre la nuit, il contient donc l'écart d'ouverture, et la dérive ajoute $\mu^2$ à son moment d'ordre deux.
\textbf{12.} Des plus hauts au cours vendeur et des plus bas au cours acheteur élargissent l'étendue et biaisent à la hausse tous les \omterm{def:rs:price-and-volume-features:range}{estimateurs de volatilité par l'étendue}, davantage pour des actions d'apparence liquide mais
à large fourchette. \textbf{13.} Yang--Zhang. \textbf{14.} Rogers--Satchell (ou Yang--Zhang avec sa partie overnight proche de zéro).
\textbf{15.} Close-to-close (une moyenne mobile exponentielle des rendements au carré), avec cinq à huit fois plus de jours pour la même précision. \textbf{16.} Que
les plus hauts et les plus bas excluent les transactions erronées et celles d'enchère, que les barres couvrent toute la séance, et quelle est la largeur de la fourchette.
\textbf{17.} Une volatilité moins bruitée fait un meilleur prédicteur de volatilité et une meilleure normalisation par le risque des autres prédicteurs. \textbf{18.}
\emph{Résultat nommé~:} par rapport aux rendements close-to-close au carré, Parkinson est 5.2 fois plus efficace, Garman--Klass 7.6 et
Rogers--Satchell 5.7, et Yang--Zhang sur vingt jours 7.5 fois la variance close-to-close sur vingt jours~; avec un écart overnight portant 30~\%
de la variance, les \omterm{def:rs:price-and-volume-features:range}{estimateurs de volatilité par l'étendue} sur un seul jour n'en voient que 63--64~\%. \textbf{19.} Une faible volatilité et l'\omterm{def:rs:price-and-volume-features:volume}{illiquidité d'Amihud} (via les
rendements en valeur absolue), et toute caractéristique normalisée par la volatilité. \textbf{20.} Le chemin~: la distance parcourue par le cours pendant la séance, et non seulement là où il
a fini.
\end{solution}

\begin{solution}{iq:rs:price-and-volume-features:1}
Parce que le rendement du mois le plus récent tend à se retourner (une autocorrélation d'ordre un négative des rendements mensuels), il joue donc contre
la continuation sur laquelle parie le momentum~; en excluant ce mois, on retire une composante de signe contraire.
Sur un marché sans ce retournement, la \omterm{def:rs:price-and-volume-features:return}{période d'exclusion} coûterait de l'information.
\iqlookfor{le retournement comme raison, et non «~convention~».}
\end{solution}

\begin{solution}{iq:rs:price-and-volume-features:2}
Utiliser l'étendue~: Garman--Klass ou, en présence d'écarts overnight, Yang--Zhang, sur le mois, après avoir nettoyé les plus hauts et les plus bas des transactions erronées~; chacun est cinq
à huit fois plus efficace que les rendements close-to-close au carré, et Yang--Zhang gère la dérive et les écarts. Corriger le biais à la baisse
d'une étendue observée en temps discret si les barres sont peu fournies.
\iqlookfor{\omterm{def:rs:price-and-volume-features:range}{estimateurs de volatilité par l'étendue}, leur efficacité et leurs cas de défaillance.}
\end{solution}

\begin{solution}{iq:rs:price-and-volume-features:3}
Le logarithme d'une moyenne de cours est, au premier ordre, une moyenne de logarithmes de cours, et chaque logarithme de cours est celui d'aujourd'hui moins une somme de rendements récents.
La différence des deux moyennes donne $w_j = (m - 1 - j)/m - (n - 1 - j)^+/n$~: des poids triangulaires croissant sur la fenêtre courte et
décroissant jusqu'à zéro sur la fenêtre longue.
\iqlookfor{la dérivation au premier ordre et la forme des poids.}
\end{solution}

\begin{solution}{iq:rs:price-and-volume-features:4}
Un test de perturbation~: recalculer chaque caractéristique après avoir multiplié toutes les entrées postérieures à une date de coupure par des facteurs aléatoires, et exiger des sorties identiques
jusqu'à la date de coupure, pour de nombreuses dates de coupure~; plus une caractéristique qui regarde volontairement vers l'avant, pour prouver que le test la détecte. Aussi une couche d'accès aux données qui
refuse les données postérieures à l'instant de décision.
\iqlookfor{un test automatisé et réfutable.}
\end{solution}

\begin{solution}{iq:rs:price-and-volume-features:5}
La variation de cours par unité de monnaie échangée, une approximation de l'impact sur le prix. Faiblesses~: elle est bruitée jour après jour (un seul mouvement important domine), elle mélange
les mouvements dictés par l'actualité à l'impact, elle dépend du niveau de prix à travers la valeur échangée, et elle vaut zéro les jours sans transaction.
\iqlookfor{ce qu'elle approxime et pourquoi elle est bruitée.}
\end{solution}

\begin{solution}{iq:rs:price-and-volume-features:6}
La caractéristique peut porter une information réelle que le simulateur ne contient pas~: un simulateur ne reproduit que les effets qu'on y a introduits. Rien
n'est appris sur la caractéristique à partir du simulateur~; le résultat sur données réelles doit être testé comme tout autre (recherche, chevauchement, coûts), et la
portée du simulateur doit être précisée~: il ne peut pas valider l'effet de cette caractéristique.
\iqlookfor{la portée d'un simulateur, et ne pas le traiter comme une preuve contre un effet réel.}
\end{solution}
